
\documentclass[10pt]{article}
\usepackage[utf8]{inputenc}
\usepackage[T1]{fontenc}
\usepackage{amsmath, amssymb, xcolor, geometry, enumitem, multicol, tcolorbox}

\definecolor{mainblue}{RGB}{0, 102, 204}
\geometry{a4paper, margin=1.2cm, top=1cm, bottom=3cm, footskip=1.2cm}

\begin{document}
\enlargethispage{2.5cm}

% Modern Header
\begin{tcolorbox}[colback=mainblue!5, colframe=mainblue, sharp corners, boxrule=0.5pt]
    \small \textbf{Unit:} undefined \hfill \textbf{Name:} \rule{3cm}{0.4pt} \\
    \small \textbf{Topic:} Variables and Expressions / Order of Operations \hfill \textbf{Date:} \rule{2cm}{0.4pt} \\
    \centering \textbf{\large Homework 2}
\end{tcolorbox}

\vspace{0.2cm}

% MCQ Section
\begin{multicols}{2}
\begin{enumerate}[label=\textbf{Q\arabic*.}, leftmargin=*, itemsep=6pt, parsep=0pt]

  \item \small What is the variable in the expression $2x + 3$?
  \begin{enumerate}[label=(\Alph*), leftmargin=0.6cm, nosep, itemsep=2pt]
    \item 2
    \item 3
    \item x
    \item 2x
    \item None of the above
  \end{enumerate}

  \item \small Identify the coefficient in the term $4y$
  \begin{enumerate}[label=(\Alph*), leftmargin=0.6cm, nosep, itemsep=2pt]
    \item y
    \item 4
    \item 4y
    \item None of the above
    \item 0
  \end{enumerate}

  \item \small What is the constant term in the expression $x + 2$?
  \begin{enumerate}[label=(\Alph*), leftmargin=0.6cm, nosep, itemsep=2pt]
    \item x
    \item 2
    \item x + 2
    \item None of the above
    \item 0
  \end{enumerate}

  \item \small Identify the operator in $5 - 2$
  \begin{enumerate}[label=(\Alph*), leftmargin=0.6cm, nosep, itemsep=2pt]
    \item 5
    \item 2
    \item -
    \item None of the above
    \item +
  \end{enumerate}

  \item \small What is the term $3x$ an example of?
  \begin{enumerate}[label=(\Alph*), leftmargin=0.6cm, nosep, itemsep=2pt]
    \item Constant
    \item Variable
    \item Coefficient times a variable
    \item Expression
    \item Equation
  \end{enumerate}

  \item \small In the expression $2x + 5$, what operation is performed first?
  \begin{enumerate}[label=(\Alph*), leftmargin=0.6cm, nosep, itemsep=2pt]
    \item Addition
    \item Multiplication
    \item Subtraction
    \item Division
    \item None of the above
  \end{enumerate}

  \item \small What does the expression $2(x + 3)$ mean?
  \begin{enumerate}[label=(\Alph*), leftmargin=0.6cm, nosep, itemsep=2pt]
    \item 2 times x plus 3
    \item 2 times the sum of x and 3
    \item x plus 3 times 2
    \item x plus 2 times 3
    \item 2x + 6
  \end{enumerate}

  \item \small Identify the expression $x - 2$ as an example of?
  \begin{enumerate}[label=(\Alph*), leftmargin=0.6cm, nosep, itemsep=2pt]
    \item A single term
    \item A constant
    \item A single-variable expression
    \item A multi-variable expression
    \item An equation
  \end{enumerate}

\end{enumerate}
\end{multicols}

\vspace{-0.2cm}
\hrule
\vspace{0.2cm}
\textbf{\small Open Ended Questions (FRQ)}
\begin{enumerate}[label=\textbf{Q\arabic*.}, start=9, itemsep=5pt, topsep=0pt]

  \item \small Explain, using examples, the difference between a variable and a constant in an algebraic expression. Provide a simple expression that includes both. \vspace{2cm}

  \item \small Draw a simple diagram to illustrate the concept of the order of operations. Use the expression $3 	imes 2 + 12 div 4$ as an example and explain each step. \vspace{2cm}

\end{enumerate}

\vfill

% Chef's Tips Footer
\begin{tcolorbox}[colback=blue!5, colframe=mainblue!50, title=\small \textbf{Chef's Tips}, fonttitle=\bfseries, bottom=1mm]
\begin{itemize}[noitemsep, topsep=2pt, leftmargin=0.5cm]
  \item \scriptsize Focus on defining key terms such as 'variable', 'constant', and 'coefficient' to ensure students have a solid foundation.\item \scriptsize Use simple expressions to illustrate the order of operations, emphasizing the sequence of steps to follow.\item \scriptsize Encourage students to use diagrams or visual aids to help them understand and explain complex concepts in a simple manner.
\end{itemize}
\end{tcolorbox}

\end{document}
  